\documentclass[11pt,a4paper]{article}
\usepackage[margin=0.85in,top=1in,bottom=1in]{geometry}
\usepackage{amsmath,amssymb,amsfonts}
\usepackage{graphicx}
\usepackage{float}
\usepackage{enumitem}
\usepackage{microtype}
\usepackage{caption}
\usepackage[hidelinks]{hyperref}
\usepackage[most,breakable]{tcolorbox}
\tcbuselibrary{breakable,skins}
\usepackage{fontspec}
\setmainfont{DejaVu Serif}
\setsansfont{DejaVu Sans}
\setmonofont{DejaVu Sans Mono}
\captionsetup{font=small,labelfont=bf,skip=4pt}
\setlist{leftmargin=1.5em,topsep=2pt}

%% Breakable problem card — prevents content cutoff across pages
\newtcolorbox{solcard}[3]{
  breakable, enhanced,
  colback=white, colframe=gray!50, arc=2pt,
  title={\textbf{#1}\hfill\textsf{\small #3\,pts}},
  fonttitle=\normalsize\sffamily\bfseries,
  before upper={\small\textit{#2}\par\smallskip\hrule height 0.4pt\smallskip},
  top=4pt, bottom=6pt, left=7pt, right=7pt,
  parbox=false,
}

\title{\textbf{Spinning top research --- Solution}}
\author{\normalsize X25 (2025) — English translation}
\date{}
\begin{document}\maketitle\vspace{-0.5em}
\begin{solcard}{A1}{Write down in vector form the expression for the instantaneous speed $\vec{v}$ of fragment $dm$ at the moment under consideration. Express your answer using $\vec{\Omega}$, $\vec{r}$.}{0.20}

We use Euler's formula for a rigid body.

\par\smallskip\noindent\textbf{Answer:} \textbackslash{}textbf{Answer:} $$ \vec{v} = [\vec{\Omega} \times  \vec{r}] $$
\end{solcard}

\begin{solcard}{A2}{Write down in vector form the expression for the angular momentum $d\vec{L}$ of the fragment $dm$ at the moment under consideration. Express your answer using $\vec{\Omega}$, $\vec{r}$, $dm$.}{0.20}

By definition, the angular momentum of an element of mass $dm$:

\par\smallskip\noindent\textbf{Answer:} \textbackslash{}textbf{Answer:} $$ d\vec{L} =dm [\vec{r}  \times  [\vec{\Omega} \times \vec{r}]] $$
\end{solcard}

\begin{solcard}{A3}{Obtain expressions for the components of the vector $d \vec{L} = (dL_x, dL_y, dL_z)$ in the coordinate system $OXYZ$. Express your answer using $X$, $Y$, $Z$, $\Omega_x$, $\Omega_y$, $\Omega_z$, $dm$.}{0.30}

Let's expand the double vector product in the expression for \(d\vec{L}\): $$ d\vec{L} =dm (\vec{r}  \times  (\vec{\Omega} \times \vec{r})) = dm (\vec{\Omega}r^2 - \vec{r}(\vec{\Omega} \cdot \vec{r})). $$Распишем поrayенное vectorное expression покомпонентно: $$ dL_x = dm\left(\Omega_x(X^2 + Y^2 + Z^2) - X(\Omega_xX + \Omega_y Y + \Omega_z Z)\right), \\ dL_y = dm\left(\Omega_y(X^2 + Y^2 + Z^2) - Y(\Omega_xX + \Omega_y Y + \Omega_z Z)\right), \\ dL_z = dm\left(\Omega_z(X^2 + Y^2 + Z^2) - Z(\Omega_xX + \Omega_y Y + \Omega_z Z)\right). $$ Simplifying, we get the answer.

\par\smallskip\noindent\textbf{Answer:} \textbackslash{}textbf{Answer:} $$ dL_x = dm\left(\Omega_x(Y^2 + Z^2) - \Omega_y XY - \Omega_z XZ \right)  \\ dL_y = dm\left( - \Omega_x YX + \Omega_y(X^2 + Z^2) - \Omega_z YZ \right)  \\ dL_z = dm\left( - \Omega_x ZX - \Omega_y ZY + \Omega_z(X^2 + Y^2)\right) $$
\end{solcard}

\begin{solcard}{A4}{For arbitrary motion of a rigid body, we can write angular momentum in the form \textbackslash{}begin{align*} L_x \&= I_{xx} \textbackslash{}Omega_x+ I_{xy} \textbackslash{}Omega_y + I_{xz} \textbackslash{}Omega_z;\textbackslash{}\textbackslash{} L_y \&= I_{yx} \textbackslash{}Omega_x + I_{yy} \textbackslash{}Omega_y + I_{yz} \textbackslash{}Omega_z;\textbackslash{}\textbackslash{} L_z \&= I_{zx} \textbackslash{}Omega_x + I_{zy} \textbackslash{}Omega_y + I_{zz} \textbackslash{}Omega_z. \textbackslash{}end{align*} Obtain expressions for the coefficients $I_{xx}$, $I_{xy}$, $I_{xz}$, $I_{yx}$, $I_{yy}$, $I_{yz}$, $I_{zx}$, $I_{zy}$, $I_{zz}$  in the form of integrals over $dm$. The integral over $dm$ is the sum over infinitesimal regions of a solid body of mass $dm$. For example, the coordinates of the center of mass can be written as $$X_c = \frac{1}{m} \int X dm; \quad Y_c = \frac{1}{m} \int Y dm, \quad   Z_c = \frac{1}{m} \int Z dm. $$ In further paragraphs, the formulas obtained here are not used.}{0.50}

For example, the coordinates of the center of mass can be written as $$X_c = \frac{1}{m} \int X dm; \quad Y_c = \frac{1}{m} \int Y dm, \quad   Z_c = \frac{1}{m} \int Z dm. $$


In further paragraphs the formulas obtained here are not used.


The answer is obtained by integrating over $dm$ the expressions of the previous paragraph.

\par\smallskip\noindent\textbf{Answer:} \textbackslash{}textbf{Answer:} $$ I_{xx} = \int{dm(Y^2 + Z^2)} \qquad I_{yy} = \int{dm(X^2 + Z^2)} \qquad I_{zz} = \int{dm(X^2 + Y^2)} \\ I_{xy} = I_{yx} = -\int{dmXY} \qquad I_{yz} = I_{zy} = -\int{dmYZ} \qquad I_{zx} = I_{xz} = -\int{dmXZ} $$
\end{solcard}

\begin{solcard}{A5}{Derive expressions for $\vec{L}$ and $T$ in the principal axes. Express your answer using $I_x$, $I_y$, $I_z$, $\omega_x$, $\omega_y$, $\omega_z$, $\vec{e}_x$, $\vec{e}_y$, $\vec{e}_z$.}{0.50}

In the principal axes, the expressions for the angular momentum components $\vec{L}$ take the form $$L_x = I_x \omega_x, \qquad \qquad L_y = I_y \omega_y, \qquad \qquad L_z = I_z \omega_z.$$This gives the answer.


According to the condition $$T = \frac{1}{2} \vec\omega \cdot \vec{L}.$$Using найденное выше expression for $\vec L$, we arrive at the answer.

\par\smallskip\noindent\textbf{Answer:} \textbackslash{}textbf{Answer:} $$ \vec{L} = I_x \omega_x \vec{e}_x + I_y \omega_y \vec{e}_y +  I_z \omega_z \vec{e}_z $$
\end{solcard}

\begin{solcard}{A6}{Obtain expressions for the components of the vector $\dot{\vec{L}}$ in the moving coordinate system $\left(\dot{\vec{L}}\right)_x$, $\left(\dot{\vec{L}}\right)_y$, $\left(\dot{\vec{L}}\right)_z$. Express your answer using $I_x$, $I_y$, $I_z$, $\omega_x$, $\omega_y$, $\omega_z$, $\dot{\omega}_x$, $\dot{\omega}_y$, $\dot{\omega}_z$.}{0.50}

Let us differentiate the previously obtained expression for $\vec L$. Since the unit vectors of the coordinate system are movable, they also need to be differentiated. $$ \dot{\vec{L}} = I_x \dfrac{d}{dt}(\omega_x \vec{e}_x) + I_y \dfrac{d}{dt}(\omega_y \vec{e}_y) +  I_z \frac{d}{dt}(\omega_z \vec{e}_z) = \\ = I_x \dot{\omega}_x \vec{e}_x + I_y \dot{\omega}_y \vec{e}_y +  I_z \dot{\omega}_z \vec{e}_z +  I_x \omega_x \dot{\vec{e}}_x + I_y \omega_y \dot{\vec{e_y}} +  I_z \omega_z \dot{\vec{e}}_z = \\ = I_x \dot{\omega}_x \vec{e}_x + I_y \dot{\omega}_y \vec{e}_y +  I_z \dot{\omega}_z \vec{e}_z + (\vec\omega \times (I_x \omega_x \vec{e}_x + I_y \omega_y \vec{e}_y + I_z \omega_z \vec{e}_z)) = \\ = I_x \dot{\omega}_x \vec{e}_x + I_y \dot{\omega}_y \vec{e}_y +  I_z \dot{\omega}_z \vec{e}_z + (\vec{\omega} \times \vec{L}) $$Распишем в координатах vectorное произведение: $$ [\vec{\omega} \times \vec{L}] = \begin{vmatrix} \vec{e}_x & \vec{e}_y & \vec{e}_z \\ \omega_x & \omega_y & \omega_z \\ L_x & L_y & L_z \end{vmatrix} = \begin{matrix} \vec{e}_x(\omega_y L_z - \omega_z L_y) + \vec{e}_y(\omega_z L_x - \omega_x L_z) + \vec{e}_z(\omega_x L_y - \omega_y L_z) \end{matrix} $$Thus we get проекции $\dot {\vec L}$ на оси системы $Oxyz$.

\par\smallskip\noindent\textbf{Answer:} \textbackslash{}textbf{Answer:} $$ \left(\dot{\vec{L}}\right)_x = I_x\dot{\omega}_x + (I_z-I_y)\omega_y \omega_z \\ \left(\dot{\vec{L}}\right)_y = I_y\dot{\omega}_y + (I_x-I_z)\omega_z \omega_x \\ \left(\dot{\vec{L}}\right)_z = I_z\dot{\omega}_z + (I_y-I_x)\omega_x \omega_y $$
\end{solcard}

\begin{solcard}{B1}{Find the projections of angular velocity $\omega_x$, $\omega_y$, $\omega_z$. Express your answer using $\omega_s$, $\dot{\theta}$, $\dot{\varphi}$, $\theta$.}{0.40}

Projecting $\vec\omega_s, \vec\omega_{\theta}, \vec\omega_{\varphi}$ onto the axes $x$, $y$, $z$ and taking into account $$\vec\omega_{\theta} = -\dot\theta \vec{e}_y, \qquad \qquad \vec\omega_{\varphi} =  \dot\varphi \vec e_Z,$$ we get the answer.

\par\smallskip\noindent\textbf{Answer:} \textbackslash{}textbf{Answer:} $$ \omega_x = \dot{\varphi} \sin{\theta} \\ \omega_y = -\dot{\theta} \\ \omega_z =  \omega_s + \dot{\varphi} \cos{\theta} $$
\begin{figure}[H]
  \centering
  \includegraphics[width=0.72\textwidth]{T2_S_EN_assets/fig1.jpg}
\end{figure}
\end{solcard}

\begin{solcard}{B2}{Show that the projection $L_1$ of the angular momentum $\vec{L}$ onto the fixed vertical axis $OZ$ is constant. Show that the projection $L_2$ of the angular momentum $\vec{L}$ onto the moving axis $Oz$ is constant.}{0.40}

The moment of gravity is expressed by the formula $$\vec{M} = \vec{l} \times m \vec g,$$ from которой it follows, что $\vec{M}$ перпендикулярен плоскости $OZz$. From the equation momentов в неподвижных осях it follows $$\dot L_1 = \dot L_Z = M_Z = 0.$$ From the equation Эйлера for оси $Oz$ it follows $$I_z \dot{\omega}_z + (I_y - I_x) \omega_x \omega_y = 0.$$ В силу симметрии $I_x = I_y$, therefore $$\dot L_2 = I_z \dot\omega_z = 0.$$ Отметим, что перпендикулярности $\vec{M}$ and $Oz$ for сохранения $L_2$ is not enough.

\end{solcard}

\begin{solcard}{B3}{Using the results of step A5, express the values ​​of the projections $L_1$ and $L_2$ in terms of $\omega_s$, $\dot{\theta}$, $\dot{\varphi}$, $\theta$, $A$, $C$.}{0.60}

Let us express $L_1$ through the projections $\vec{L}$: $$ L_1 = L_z \cos{\theta} + L_x\sin{\theta} = A \omega_x \sin{\theta} + C \omega_z \cos{\theta}. $$ Using result partа B1, we get окончательный answer for $L_1$: $$ L_1 =C(\omega_s + \dot{\varphi}\cos{\theta})\cos{\theta} + A\dot{\varphi}\sin^2{\theta}. $$For $L_2$ true $$L_2 = L_z = C\omega_z = C(\omega_s + \dot{\varphi}\cos{\theta}).$$

\par\smallskip\noindent\textbf{Answer:} \textbackslash{}textbf{Answer:} $$ L_1 =C(\omega_s + \dot{\varphi}\cos{\theta})\cos{\theta} + A\dot{\varphi}\sin^2{\theta} \\ \ \\ L_2 = C(\omega_s + \dot{\varphi}\cos{\theta}) $$
\end{solcard}

\begin{solcard}{B4}{At this point we will assume that $l = 0$, and the moment of gravity can be ignored. Let the angular momentum of the top be equal to $L_0$ and directed along the axis $OZ$. At time $t = 0$ the values ​​of angles $\theta = \theta_0$, $\varphi  = \varphi_0$. Find the dependence of angles $\theta$, $\varphi$ on time. This movement is called regular precession.}{0.70}

In the described case $$L_1 = L_0, \qquad \qquad L_2 = L_0 \cos\theta.$$ В силу постоянства величин $L_0$, $L_1$, $L_2$ angle $\theta$ also остаётся постоянным. Using выражения for $L_1$, $L_2$ from предыдущего partа, we get систему уравнений: $$ BEEQUATION $$Solving the system, we get $$\dot\varphi = \frac{L_0}{A}$$Next зависимость $\varphi$ from time to time turns out to be simple integration.

\par\smallskip\noindent\textbf{Answer:} \textbackslash{}textbf{Answer:} $$ \theta(t) = \theta_0 \\ \varphi(t) = \varphi_0 + \frac{L_0}{A}t $$
\end{solcard}

\begin{solcard}{С1}{For an arbitrary position of the top, express the total energy $E$ of the motion of the top (the sum of kinetic and potential) in terms of $\omega_s$, $\dot{\theta}$, $\dot{\varphi}$, $A$, $C$, $m$, $g$, $l$, $\theta$.}{0.60}

Let's write down the expressions for potential and kinetic energies: $$T = \frac{1}{2} (C\omega_z^2 + A(\omega_x^2 + \omega_y^2) ) = \frac{1}{2} (C(\omega_s + \dot\varphi \cos\theta)^2 + A(\dot\theta^2 + \dot\varphi^2 \sin^2 \theta)), \\ \Pi = mgl \cos\theta.$$Полная energy: $$ E = T + \Pi. $$

\par\smallskip\noindent\textbf{Answer:} \textbackslash{}textbf{Answer:} $$ E = \frac{1}{2} (C(\omega_s + \dot\varphi \cos\theta)^2 + A(\dot\theta^2 + \dot\varphi^2 \sin^2 \theta)) + mgl \cos\theta $$
\end{solcard}

\begin{solcard}{С2}{Express $\omega_s$ and $\dot{\varphi}$ in terms of $L_1$, $L_2$, $A$, $C$, $ \theta$ using the results from step B3.}{0.30}

The answer is obtained by solving the system for $\omega_s$, $\dot\varphi$. \textbackslash{}begin{equation} \textbackslash{}left\textbackslash{}{ \textbackslash{}begin{aligned} L_1 \&= \textbackslash{}dot{\textbackslash{}varphi}(A\textbackslash{}sin^2{\textbackslash{}theta} + C\textbackslash{}cos^2{\textbackslash{}theta}) + C\textbackslash{}omega_s\textbackslash{}cos{\textbackslash{}theta} \textbackslash{}\textbackslash{} L_2 \&= C(\textbackslash{}omega_s + \textbackslash{}dot{\textbackslash{}varphi}\textbackslash{}cos{\textbackslash{}theta}) \textbackslash{}end{aligned} \textbackslash{}right. \textbackslash{}end{equation}

\par\smallskip\noindent\textbf{Answer:} \textbackslash{}textbf{Answer:} $$ \dot{\varphi} = \frac{L_1 - L_2\cos{\theta}}{A\sin^2{\theta}}\\ \ \\ \omega_s = \frac{L_2}{C} - \frac{L_1 - L_2\cos{\theta}}{A\sin^2{\theta}}\cos{\theta} $$
\end{solcard}

\begin{solcard}{C3}{Using the constancy $L_1$, $L_2$, reduce the expression for the total energy $E$ to the form $$E = \frac{\mu}{2} \dot{\theta}^2 + U(\theta).$$ Поrayandте выражения for $\mu$, $U(\theta)$ via/through $A$, $C$, $L_1$, $L_2$, $m$,$g$, $l$. Note. As is known, the total energy is determined up to a constant, which for convenience can be set equal to zero.}{0.60}

Note. As is known, the total energy is determined up to a constant, which for convenience can be set equal to zero.


Substitution into the expression for the energy $\omega_s$ and $\dot{\varphi}$ from the previous paragraph leads to the result $$E = \frac{A}{2} \dot\theta^2+\frac{(L_1 - L_2 \cos\theta)^2}{2 A \sin^2\theta} + mgl \cos\theta + \frac{L_2^2}{2C}.$$ Последнее константное term можно отбросить: $$E = \frac{A}{2} \dot\theta^2+\frac{(L_1 - L_2 \cos\theta)^2}{2 A \sin^2\theta} + mgl \cos\theta.$$ The resulting expression yields the answer.

\par\smallskip\noindent\textbf{Answer:} \textbackslash{}textbf{Answer:} $$ \mu = A \\ U(\theta) = mgl\cos{\theta} + \frac{(L_1 - L_2\cos{\theta})^2}{2A\sin^2{\theta}} $$
\end{solcard}

\begin{solcard}{D1}{Find the period of oscillation of the top. You will receive the answer in the form of a definite integral containing $A$, $C$, $l$, $\theta_0$, $g$. Also obtain a simplified expression for the case $\pi - \theta_0 \ll 1$.}{0.80}

For given initial conditions $$L_1 = L_2 = 0.$$ Energy conservation law в этом сrayае запишется в виде $$\frac{A\dot{\theta}^2}{2} + mgl\cos{\theta} = mgl\cos{\theta_0}.$$ Hence выражается зависимость $\dot\theta(\theta)$: $$\dot{\theta} = \pm \sqrt{\frac{2mgl}{A}(\cos{\theta_0} - \cos{\theta})}.$$ Движение от $\theta = \theta_0$ до $\theta = \pi$ с $\dot\theta > 0$ занимает четверть periodа. For нахождения этого времени требуется проинтегрировать equation $$dt = \frac{d\theta}{\dot\theta}.$$Integrating: $$\frac{T}{4} = \sqrt{\frac{A}{2mgl}}\int\limits_{\theta_0}^{\pi}{\frac{d\theta}{\sqrt{\cos{\theta_0} - \cos{\theta}}}}.$$From here follows the answer.


Let's introduce the variable $\beta = \pi - \theta$, then $\beta \ll 1$. Method 1 Taking into account the smallness of $\beta$, the expression for energy will be written as: $$ E = \frac{A\dot{\beta}^2}{2} + mgl\frac{\beta^2}{2}\\ $$ Такой вид уравнения соответствует гармоническим колебаниям с periodом $$ T = 2\pi \sqrt{\frac{A}{mgl}}. $$ 2 метод Введём $\beta_0 = \pi - \theta_0$. Transforming интеграл $$\int\limits_{\theta_0}^{\pi}{\frac{d\theta}{\sqrt{\cos{\theta_0} - \cos{\theta}}}}$$ в рассматриваемом for/atближении. $$\int\limits_{\theta_0}^{\pi}{\frac{d\theta}{\sqrt{\cos{\theta_0} - \cos{\theta}}}} = \int\limits_{\beta_0}^{0}{\frac{-d\beta}{\sqrt{\cos{(\pi - \beta_0)} - \cos{(\pi - \beta)}}}} = \int\limits_{0}^{\beta_0}{\frac{d\beta}{\sqrt{\cos \beta - \cos \beta_0}}} \simeq \\ \simeq \sqrt{2} \int\limits_{0}^{\beta_0}{\frac{d\beta}{\sqrt{\beta_0^2 - \beta^2}}} = \sqrt{2} \int\limits_{0}^{\beta_0}{\frac{\frac{d\beta}{\beta_0}}{\sqrt{1 - \frac{\beta^2}{\beta_0^2}}}} = \sqrt{2} \int\limits_{0}^{1}{\frac{dx}{\sqrt{1 -x^2}}} = \sqrt{2} \frac{\pi}{2}.$$ Подстановка найденного значения в исходное expression for $T$ leads to the answer obtained method 1.


1 method


Taking into account the smallness of $\beta$, the expression for energy will be written in the form: $$ E = \frac{A\dot{\beta}^2}{2} + mgl\frac{\beta^2}{2}\\ $$ Такой вид уравнения соответствует гармоническим колебаниям с periodом $$ T = 2\pi \sqrt{\frac{A}{mgl}}. $$


Method 2


Let's enter $\beta_0 = \pi - \theta_0$. Let's transform the integral


$$\int\limits_{\theta_0}^{\pi}{\frac{d\theta}{\sqrt{\cos{\theta_0} - \cos{\theta}}}}$$ in the considered approximation.


$$\int\limits_{\theta_0}^{\pi}{\frac{d\theta}{\sqrt{\cos{\theta_0} - \cos{\theta}}}} = \int\limits_{\beta_0}^{0}{\frac{-d\beta}{\sqrt{\cos{(\pi - \beta_0)} - \cos{(\pi - \beta)}}}} = \int\limits_{0}^{\beta_0}{\frac{d\beta}{\sqrt{\cos \beta - \cos \beta_0}}} \simeq \\ \simeq \sqrt{2} \int\limits_{0}^{\beta_0}{\frac{d\beta}{\sqrt{\beta_0^2 - \beta^2}}} = \sqrt{2} \int\limits_{0}^{\beta_0}{\frac{\frac{d\beta}{\beta_0}}{\sqrt{1 - \frac{\beta^2}{\beta_0^2}}}} = \sqrt{2} \int\limits_{0}^{1}{\frac{dx}{\sqrt{1 -x^2}}} = \sqrt{2} \frac{\pi}{2}.$$ Подстановка найденного значения в исходное expression for $T$ results in the answer obtained in method 1.

\par\smallskip\noindent\textbf{Answer:} \textbackslash{}textbf{Answer:} $$ T = 4\sqrt{\frac{A}{2mgl}}\int\limits_{\theta_0}^{\pi}{\frac{d\theta}{\sqrt{\cos{\theta_0} - \cos{\theta}}}} $$
\end{solcard}

\begin{solcard}{D2}{Express $L_1$ and $L_2$ in terms of $A$, $C$, $\omega_0$.}{0.20}

In a given situation $L_1 = L_2$. Their values ​​are instantly obtained from the expression $\vec{L}$ in the principal axes.

\par\smallskip\noindent\textbf{Answer:} \textbackslash{}textbf{Answer:} $$ L_1 = L_2 = C \omega_0 $$
\end{solcard}

\begin{solcard}{D3}{Find a condition relating $A$, $C$, $m$, $g$, $l$, $\omega_0$ such that the position $\theta = 0$ is stable for the deviations described above. This condition is called Maievsky's condition. It turns out that if the found condition is satisfied, the equilibrium position is stable even if the conservation condition $L_1$ and $L_2$ is abandoned.}{0.50}

At $L_1 = L_2 = C\omega_0$ the expression for the energy becomes $$E = \frac{A}{2} \dot\theta^2 + mgl \cos \theta + \frac{C^2 \omega_0^2 (1-\cos \theta)^2}{2A\sin^2\theta}.$$ Expanding taking into account малости $\theta$: $$E = \frac{A\dot\theta^2}{2} + \left(\frac{C^2 \omega_0^2}{4A} - mgl \right)\frac{\theta^2}{2}.$$ For устойчивости необходимо, чтобы coefficient перед $\theta^2$ was positive, which determines the required condition.

\par\smallskip\noindent\textbf{Answer:} \textbackslash{}textbf{Answer:} $$ \frac{C^2\omega_0^2}{4A} > mgl $$
\end{solcard}

\begin{solcard}{D4}{Obtain an equation for $\theta$ that is satisfied by the angle $\theta_1$. The equation may contain $A$, $C$, $L_1$, $L_2$, $m$, $g$, $l$. Clue. $U(\theta)$ can be represented as a function of the new variable $s = \cos \theta$, which greatly simplifies the solution.}{0.70}

Clue. $U(\theta)$ can be represented as a function of the new variable $s = \cos \theta$, which greatly simplifies the solution.


In the equilibrium position $\cfrac{d}{d\theta}U(\theta) = 0$. Let's move on to the variable $s = \cos \theta$. Then $$\frac{d}{d\theta} U(\theta) = \frac{d}{ds} U(s) \cdot \frac{ds}{d\theta} = \frac{d}{ds}U(s) \cdot \sin\theta.$$ hence it follows, что condition $$\frac{d}{d\theta}U(\theta) = 0$$ эквивалентно условию $$\frac{d}{ds} U(s) = 0$$ for all $\theta \in (0, \pi/2)$. We get $\cfrac{d}{ds}U(s)$ $$U(s) = mgls + \frac{\left(L_1 - L_2s\right) ^ 2}{2A\left(1-s^2\right)} \\ \frac{d}{d\theta}U(s) = mgl - \frac{(L_1 - L_2s)L_2}{A(1 - s^2)} +\frac{(L_1 - L_2s)^2}{A(1-s^2)^2}.$$Упростим expression and перейдём обратно к переменной $\theta$. Equating it to zero, we obtain the required equation.

\par\smallskip\noindent\textbf{Answer:} \textbackslash{}textbf{Answer:} $$ \frac{(L_2 \cos \theta - L_1)(L_2 - L_1 \cos \theta)}{\sin^4 \theta} + Amgl = 0 $$
\end{solcard}

\begin{solcard}{D5}{Find the possible precession rates of the top $\omega_p$ at $\theta = \theta_1$. Express your answer using $A$, $C$, $m$, $g$, $l$, $\omega_s$, $\theta_1$. One of the solutions can be to substitute expressions for $L_1$, $L_2$ into the equation from D4 through the angular velocities $\omega_s$, $\omega_p$.}{0.80}

Substituting the expressions for $L_1$ and $L_2$ from point B3 into the equation obtained in the previous point, we arrive at a quadratic equation for $\omega_p$: $$\omega_p^2 (A-C) \cos\theta_1 - \omega_p C\omega_s + mgl = 0.$$ It has two solutions, which is the answer.

\par\smallskip\noindent\textbf{Answer:} \textbackslash{}textbf{Answer:} $$ \omega_p = \frac{C \omega_s}{2(A-C) \cos \theta_1}\left(1 \pm \sqrt{1 - \frac{4mgl(A-C) \cos\theta_1}{C^2 \omega_s^2}}\right) $$
\end{solcard}

\begin{solcard}{D6}{Obtain approximate expressions for the frequencies in the case of $mgl \ll C \omega_s^2$. Consider that $A$ and $C$ are quantities of the same order.}{0.30}

Let us expand $\omega_p$ using the approximate formula $\sqrt{1 + x} \simeq 1 + \frac{x}{2}$ for $x \ll 1$. For a larger root:


For a smaller root:


Note that the frequency $\omega_p^{slow}$ coincides with the precession frequency of the “fast” top in the simplified model, that is, when the contribution of precession to the angular momentum is neglected.

\par\smallskip\noindent\textbf{Answer:} \textbackslash{}textbf{Answer:} $$ \omega_p^{fast} = \frac{C\omega_s}{(A-C)\cos\theta_1} $$
\end{solcard}

\begin{solcard}{D7}{Determine the angle $\alpha_{max}$ of the maximum deviation of the top axis from the horizontal during the subsequent movement. Express your answer in terms of $A$, $C$, $m$, $g$, $l$, $\omega_0$.}{0.70}

Let us write the expression for energy taking into account the initial conditions: $$ L_1 = 0, \\ L_2 = C\omega_0. \\ E =\frac{A\dot{\theta}^2}{2} + mgl\cos{\theta} + \frac{C^2 \omega_0^2 \cos^2{\theta}}{2A\sin^2{\theta}} = 0. $$ Перейдём к переменной $\alpha$ по формуле: $$ \theta = \frac{\pi }{2} + \alpha. $$ Then $$ \frac{A\dot{\alpha}^2}{2}  -mgl\sin{\alpha} + \frac{C^2 \omega_0^2 \sin^2{\alpha}}{2A\cos^2{\alpha}} = 0. $$Condition $\alpha = \alpha_{max}$ соответствует $\dot{\alpha} = 0$. Substituting $\dot\alpha = 0$ в equation выше and преобсовывая, we get квадратное equation на синус максимального угла отклонения: $$ \sin^2{\alpha_{max}} + \frac{C^2 \omega_0^2}{2Amgl}\sin{\alpha_{max}} - 1 = 0. $$ Решим поrayенное квадратное equation: $$ \sin{\alpha_{max}} = \pm\sqrt{1+\left(\frac{C^2 \omega_0^2}{4Amgl}\right)^2}-\frac{C^2 \omega_0^2}{4Amgl} $$ Since $|\sin \alpha_{max}| \leq 1$, it follows выбрать корень с «$+$".

\par\smallskip\noindent\textbf{Answer:} \textbackslash{}textbf{Answer:} $$ \alpha_{max} = \arcsin \left(\sqrt{1+\left(\frac{C^2 \omega_0^2}{4Amgl}\right)^2}-\frac{C^2 \omega_0^2}{4Amgl}\right) $$
\end{solcard}

\begin{solcard}{D8}{Find the dependence of the angle of deviation of the top axis from the horizontal on time $\alpha(t)$. Express your answer in terms of $A$, $C$, $m$, $g$, $l$, $\omega_0$.}{0.70}

Let us write an approximate expression for $\alpha_{max}$ in this case: $$ \alpha_{max} = \frac{2Amgl}{C^2\omega^2_0} \sim \frac{mgl}{C \omega_0^2} \ll 1. $$ В силу малости $\alpha_{max}$ можем рассматривать expression for энергии for малых $\alpha$. Taking into account начальных условий полная energy equals нулю. $$ \frac{A\dot{\alpha}^2}{2} + \frac{C^2 \omega_0^2}{2A} \alpha^2  -mgl\alpha = 0 $$ Такой вид выражения for энергии соответствует колебаниям с частотой $$\omega = \frac{C}{A} \omega_0,$$ амплитудой $$\alpha_0 = \frac{Amgl}{C^2 \omega_0^2}$$ and смещением от центра $$\alpha_1 = \alpha_0 = \frac{Amgl}{C^2 \omega_0^2}.$$Taking into account всех найденных parameterов движения and начальных условий we get зависимость $\alpha(t)$.

\par\smallskip\noindent\textbf{Answer:} \textbackslash{}textbf{Answer:} $$ \alpha(t) = \frac{Amgl}{C^2 \omega_0^2} \left(1 - \cos \left(\frac{C}{A}\omega_0 t\right)\right) $$
\end{solcard}

\begin{solcard}{D9}{Find the dependence of the precession angle on time $\varphi(t)$. At the initial moment $\varphi(0) = \varphi_0$. Express your answer using $A$, $C$, $m$, $g$, $l$, $\omega_0$, $\varphi_0$.}{0.50}

Let's express $\dot{\varphi}(t)$ using the found dependence $\alpha(t)$: $$ \dot{\varphi} = \frac{L_1 - L_2\cos{\theta}}{A\sin^2{\theta}} = \frac{C\omega_0\sin{\alpha}}{A\cos^2{\alpha}} \simeq \frac{mgl}{C\omega_0}\left(1 - \cos \left(\frac{C}{A}\omega_0 t\right)\right) $$ Integrating, we get the answer.

\par\smallskip\noindent\textbf{Answer:} \textbackslash{}textbf{Answer:} $$ \varphi(t) = \varphi_0 + \frac{mgl}{C\omega_0} \left(t - \frac{A}{C\omega_0} \sin \left(\frac{C}{A} \omega_0 t\right)\right) $$
\end{solcard}

\begin{solcard}{D10}{Draw a characteristic graph of the dependence $\alpha(\varphi)$. Determine the angle $\Delta \varphi$ through which the top axis rotates between two successive returns of the top axis to a horizontal position. Express your answer using $A$, $C$, $m$, $g$, $l$, $\omega_0$.}{0.50}

Expressions $\varphi(t)$, $\alpha(t)$ parametrically define a cycloid, the characteristic form of which is presented in the figure below.


The time between subsequent instants $\alpha = 0$ is $t = \cfrac{2 \pi}{\omega_0}\cfrac{A}{C}$. Substituting $t$ into the expression for $\varphi$ gives the answer.

\begin{figure}[H]
  \centering
  \includegraphics[width=0.72\textwidth]{T2_S_EN_assets/fig2.jpg}
\end{figure}
\end{solcard}

\begin{solcard}{D11}{Calculate $\Delta \varphi$ for the described system.}{0.50}

Let's find $A$, $C$, $ l$: $$ C = \frac{(2m)R^2}{2} =  mR^2 \\ A = \frac{(2m)R^2}{4} + (2m)R^2 + \frac{mR^2}{3} = \frac{17}{6}mR^2 \\ l = \frac{2m  R + m \frac{R}{2}}{3m} = \frac{5}{6}R $$Substituting в формулу for $\Delta \varphi$ найденные величины (требуется also учесть, что mass волчка составляет $3m$) and get the answer.

\par\smallskip\noindent\textbf{Answer:} \textbackslash{}textbf{Answer:} $$ \Delta \varphi = \frac{17}{120} \pi $$
\end{solcard}

\end{document}